% Seminar 3 - Pagina 4. Fragment pentru
% % Seminar 3 - Pagina 4. Fragment pentru
% % Seminar 3 - Pagina 4. Fragment pentru
% % Seminar 3 - Pagina 4. Fragment pentru
% \input{seminars/pages/seminar3-pagina4-ro}; fără preambul sau pagină nouă.
% NOTE EDITORIALE:
% - Primele trei rânduri continuă exercițiul 2 de pe pagina 3.
% - Cardinalul intersecției im(f) cu [n,n+1) pare să fie <= 1.
% - Simbolurile pentru relația de echivalență și sistemul de reprezentanți
%   au fost culese tipografic pentru lizibilitate, fără a schimba demonstrația.
% - Manuscrisul scrie S sub familia (x_i)_{i in I} și rezumă condițiile.
% - Diagrama proprietății de universalitate este redesenată sub forma
%   unui triunghi comutativ, cu aceleași aplicații.
% - Teorema proprietății de universalitate este doar enunțată aici.

\subsection*{Seminar 3 - Pagina 4}
\begingroup
\setlength{\abovedisplayskip}{5pt}
\setlength{\belowdisplayskip}{5pt}
\setlength{\abovedisplayshortskip}{3pt}
\setlength{\belowdisplayshortskip}{3pt}

\noindent\textbf{Exercițiul 2 (continuare).} Din pagina anterioară,
\[
 \operatorname{im}(f)=\bigcup_{n\in\mathbb Z}
 \bigl(\operatorname{im}(f)\cap[n,n+1)\bigr).
\]
Din $|f(x)-f(y)|>1$ pentru $x\neq y$ rezultă
$|\operatorname{im}(f)\cap[n,n+1)|\leq1$.
Deci $\operatorname{im}(f)$ este finită sau numărabilă: contradicție.

\medskip
\noindent\underline{Relații de echivalență}

\noindent O relație $R\subseteq M\times M$ se numește
\emph{relație de echivalență} dacă satisface:
\[
\begin{array}{ll}
 \text{(R) Reflexivitate:}&(x,x)\in R,\\[2pt]
 \text{(S) Simetrie:}&(x,y)\in R\ \Longrightarrow\ (y,x)\in R,\\[2pt]
 \text{(T) Tranzitivitate:}&
  \left.\begin{array}{l}(x,y)\in R\\(y,z)\in R\end{array}\right\}
  \Longrightarrow (x,z)\in R.
\end{array}
\]
Notăm $x\sim y$ dacă $(x,y)\in R$. Clasa de echivalență a
lui $x\in M$ este
\[
 \widehat{x}=\{y\in M\mid y\sim x\}.
\]
Mulțimea factor și proiecția canonică sunt
\[
 M/{\sim}:=\{\widehat{x}\mid x\in M\},
 \qquad
 \pi:M\longrightarrow M/{\sim},
 \quad \pi(x)=\widehat{x}.
\]

\noindent Un \emph{sistem de reprezentanți} este o familie
$S=(x_i)_{i\in I}$ astfel încât
\[
 \begin{cases}
   x_i\not\sim x_j,&i\neq j,\\[2pt]
   \forall x\in M\ \exists i\in I:\ x\sim x_i.
 \end{cases}
\]
Pentru submulțimea corespunzătoare $S\subseteq M$, restricția
$\pi|_S:S\to M/{\sim}$ este bijectivă.

\medskip
\noindent\underline{Proprietatea de universalitate}

\noindent\textbf{Observație.} Fie $f:M\to X$. Definim
\[
 R_f:=\{(x,y)\in M\times M\mid f(x)=f(y)\}.
\]
Atunci $x\sim y\iff f(x)=f(y)$ definește relația de echivalență
indusă de $f$ pe $M$.

\noindent\fbox{T}\quad \textbf{Proprietatea de universalitate a mulțimii factor.}
Fie $\sim$ o relație de echivalență pe $M$ și fie
$\pi:M\to M/{\sim}$ proiecția canonică.
Pentru orice aplicație $f:M\to X$ cu $\sim\,\subseteq R_f$,
există o unică aplicație $\overline f:M/{\sim}\to X$ astfel încât
\[
 \overline f\circ\pi=f.
\]
\[
\begin{tikzcd}[column sep=large,row sep=large]
 M\arrow[r,"\pi"]\arrow[dr,"f"']
   & M/{\sim}\arrow[d,dashed,"{\overline f}"]\\
 & X
\end{tikzcd}
\]
% Manuscrisul se oprește după enunțul proprietății de universalitate.

\endgroup
; fără preambul sau pagină nouă.
% NOTE EDITORIALE:
% - Primele trei rânduri continuă exercițiul 2 de pe pagina 3.
% - Cardinalul intersecției im(f) cu [n,n+1) pare să fie <= 1.
% - Simbolurile pentru relația de echivalență și sistemul de reprezentanți
%   au fost culese tipografic pentru lizibilitate, fără a schimba demonstrația.
% - Manuscrisul scrie S sub familia (x_i)_{i in I} și rezumă condițiile.
% - Diagrama proprietății de universalitate este redesenată sub forma
%   unui triunghi comutativ, cu aceleași aplicații.
% - Teorema proprietății de universalitate este doar enunțată aici.

\subsection*{Seminar 3 - Pagina 4}
\begingroup
\setlength{\abovedisplayskip}{5pt}
\setlength{\belowdisplayskip}{5pt}
\setlength{\abovedisplayshortskip}{3pt}
\setlength{\belowdisplayshortskip}{3pt}

\noindent\textbf{Exercițiul 2 (continuare).} Din pagina anterioară,
\[
 \operatorname{im}(f)=\bigcup_{n\in\mathbb Z}
 \bigl(\operatorname{im}(f)\cap[n,n+1)\bigr).
\]
Din $|f(x)-f(y)|>1$ pentru $x\neq y$ rezultă
$|\operatorname{im}(f)\cap[n,n+1)|\leq1$.
Deci $\operatorname{im}(f)$ este finită sau numărabilă: contradicție.

\medskip
\noindent\underline{Relații de echivalență}

\noindent O relație $R\subseteq M\times M$ se numește
\emph{relație de echivalență} dacă satisface:
\[
\begin{array}{ll}
 \text{(R) Reflexivitate:}&(x,x)\in R,\\[2pt]
 \text{(S) Simetrie:}&(x,y)\in R\ \Longrightarrow\ (y,x)\in R,\\[2pt]
 \text{(T) Tranzitivitate:}&
  \left.\begin{array}{l}(x,y)\in R\\(y,z)\in R\end{array}\right\}
  \Longrightarrow (x,z)\in R.
\end{array}
\]
Notăm $x\sim y$ dacă $(x,y)\in R$. Clasa de echivalență a
lui $x\in M$ este
\[
 \widehat{x}=\{y\in M\mid y\sim x\}.
\]
Mulțimea factor și proiecția canonică sunt
\[
 M/{\sim}:=\{\widehat{x}\mid x\in M\},
 \qquad
 \pi:M\longrightarrow M/{\sim},
 \quad \pi(x)=\widehat{x}.
\]

\noindent Un \emph{sistem de reprezentanți} este o familie
$S=(x_i)_{i\in I}$ astfel încât
\[
 \begin{cases}
   x_i\not\sim x_j,&i\neq j,\\[2pt]
   \forall x\in M\ \exists i\in I:\ x\sim x_i.
 \end{cases}
\]
Pentru submulțimea corespunzătoare $S\subseteq M$, restricția
$\pi|_S:S\to M/{\sim}$ este bijectivă.

\medskip
\noindent\underline{Proprietatea de universalitate}

\noindent\textbf{Observație.} Fie $f:M\to X$. Definim
\[
 R_f:=\{(x,y)\in M\times M\mid f(x)=f(y)\}.
\]
Atunci $x\sim y\iff f(x)=f(y)$ definește relația de echivalență
indusă de $f$ pe $M$.

\noindent\fbox{T}\quad \textbf{Proprietatea de universalitate a mulțimii factor.}
Fie $\sim$ o relație de echivalență pe $M$ și fie
$\pi:M\to M/{\sim}$ proiecția canonică.
Pentru orice aplicație $f:M\to X$ cu $\sim\,\subseteq R_f$,
există o unică aplicație $\overline f:M/{\sim}\to X$ astfel încât
\[
 \overline f\circ\pi=f.
\]
\[
\begin{tikzcd}[column sep=large,row sep=large]
 M\arrow[r,"\pi"]\arrow[dr,"f"']
   & M/{\sim}\arrow[d,dashed,"{\overline f}"]\\
 & X
\end{tikzcd}
\]
% Manuscrisul se oprește după enunțul proprietății de universalitate.

\endgroup
; fără preambul sau pagină nouă.
% NOTE EDITORIALE:
% - Primele trei rânduri continuă exercițiul 2 de pe pagina 3.
% - Cardinalul intersecției im(f) cu [n,n+1) pare să fie <= 1.
% - Simbolurile pentru relația de echivalență și sistemul de reprezentanți
%   au fost culese tipografic pentru lizibilitate, fără a schimba demonstrația.
% - Manuscrisul scrie S sub familia (x_i)_{i in I} și rezumă condițiile.
% - Diagrama proprietății de universalitate este redesenată sub forma
%   unui triunghi comutativ, cu aceleași aplicații.
% - Teorema proprietății de universalitate este doar enunțată aici.

\subsection*{Seminar 3 - Pagina 4}
\begingroup
\setlength{\abovedisplayskip}{5pt}
\setlength{\belowdisplayskip}{5pt}
\setlength{\abovedisplayshortskip}{3pt}
\setlength{\belowdisplayshortskip}{3pt}

\noindent\textbf{Exercițiul 2 (continuare).} Din pagina anterioară,
\[
 \operatorname{im}(f)=\bigcup_{n\in\mathbb Z}
 \bigl(\operatorname{im}(f)\cap[n,n+1)\bigr).
\]
Din $|f(x)-f(y)|>1$ pentru $x\neq y$ rezultă
$|\operatorname{im}(f)\cap[n,n+1)|\leq1$.
Deci $\operatorname{im}(f)$ este finită sau numărabilă: contradicție.

\medskip
\noindent\underline{Relații de echivalență}

\noindent O relație $R\subseteq M\times M$ se numește
\emph{relație de echivalență} dacă satisface:
\[
\begin{array}{ll}
 \text{(R) Reflexivitate:}&(x,x)\in R,\\[2pt]
 \text{(S) Simetrie:}&(x,y)\in R\ \Longrightarrow\ (y,x)\in R,\\[2pt]
 \text{(T) Tranzitivitate:}&
  \left.\begin{array}{l}(x,y)\in R\\(y,z)\in R\end{array}\right\}
  \Longrightarrow (x,z)\in R.
\end{array}
\]
Notăm $x\sim y$ dacă $(x,y)\in R$. Clasa de echivalență a
lui $x\in M$ este
\[
 \widehat{x}=\{y\in M\mid y\sim x\}.
\]
Mulțimea factor și proiecția canonică sunt
\[
 M/{\sim}:=\{\widehat{x}\mid x\in M\},
 \qquad
 \pi:M\longrightarrow M/{\sim},
 \quad \pi(x)=\widehat{x}.
\]

\noindent Un \emph{sistem de reprezentanți} este o familie
$S=(x_i)_{i\in I}$ astfel încât
\[
 \begin{cases}
   x_i\not\sim x_j,&i\neq j,\\[2pt]
   \forall x\in M\ \exists i\in I:\ x\sim x_i.
 \end{cases}
\]
Pentru submulțimea corespunzătoare $S\subseteq M$, restricția
$\pi|_S:S\to M/{\sim}$ este bijectivă.

\medskip
\noindent\underline{Proprietatea de universalitate}

\noindent\textbf{Observație.} Fie $f:M\to X$. Definim
\[
 R_f:=\{(x,y)\in M\times M\mid f(x)=f(y)\}.
\]
Atunci $x\sim y\iff f(x)=f(y)$ definește relația de echivalență
indusă de $f$ pe $M$.

\noindent\fbox{T}\quad \textbf{Proprietatea de universalitate a mulțimii factor.}
Fie $\sim$ o relație de echivalență pe $M$ și fie
$\pi:M\to M/{\sim}$ proiecția canonică.
Pentru orice aplicație $f:M\to X$ cu $\sim\,\subseteq R_f$,
există o unică aplicație $\overline f:M/{\sim}\to X$ astfel încât
\[
 \overline f\circ\pi=f.
\]
\[
\begin{tikzcd}[column sep=large,row sep=large]
 M\arrow[r,"\pi"]\arrow[dr,"f"']
   & M/{\sim}\arrow[d,dashed,"{\overline f}"]\\
 & X
\end{tikzcd}
\]
% Manuscrisul se oprește după enunțul proprietății de universalitate.

\endgroup
; fără preambul sau pagină nouă.
% NOTE EDITORIALE:
% - Primele trei rânduri continuă exercițiul 2 de pe pagina 3.
% - Cardinalul intersecției im(f) cu [n,n+1) pare să fie <= 1.
% - Simbolurile pentru relația de echivalență și sistemul de reprezentanți
%   au fost culese tipografic pentru lizibilitate, fără a schimba demonstrația.
% - Manuscrisul scrie S sub familia (x_i)_{i in I} și rezumă condițiile.
% - Diagrama proprietății de universalitate este redesenată sub forma
%   unui triunghi comutativ, cu aceleași aplicații.
% - Teorema proprietății de universalitate este doar enunțată aici.

\subsection*{Seminar 3 - Pagina 4}
\begingroup
\setlength{\abovedisplayskip}{5pt}
\setlength{\belowdisplayskip}{5pt}
\setlength{\abovedisplayshortskip}{3pt}
\setlength{\belowdisplayshortskip}{3pt}

\noindent\textbf{Exercițiul 2 (continuare).} Din pagina anterioară,
\[
 \operatorname{im}(f)=\bigcup_{n\in\mathbb Z}
 \bigl(\operatorname{im}(f)\cap[n,n+1)\bigr).
\]
Din $|f(x)-f(y)|>1$ pentru $x\neq y$ rezultă
$|\operatorname{im}(f)\cap[n,n+1)|\leq1$.
Deci $\operatorname{im}(f)$ este finită sau numărabilă: contradicție.

\medskip
\noindent\underline{Relații de echivalență}

\noindent O relație $R\subseteq M\times M$ se numește
\emph{relație de echivalență} dacă satisface:
\[
\begin{array}{ll}
 \text{(R) Reflexivitate:}&(x,x)\in R,\\[2pt]
 \text{(S) Simetrie:}&(x,y)\in R\ \Longrightarrow\ (y,x)\in R,\\[2pt]
 \text{(T) Tranzitivitate:}&
  \left.\begin{array}{l}(x,y)\in R\\(y,z)\in R\end{array}\right\}
  \Longrightarrow (x,z)\in R.
\end{array}
\]
Notăm $x\sim y$ dacă $(x,y)\in R$. Clasa de echivalență a
lui $x\in M$ este
\[
 \widehat{x}=\{y\in M\mid y\sim x\}.
\]
Mulțimea factor și proiecția canonică sunt
\[
 M/{\sim}:=\{\widehat{x}\mid x\in M\},
 \qquad
 \pi:M\longrightarrow M/{\sim},
 \quad \pi(x)=\widehat{x}.
\]

\noindent Un \emph{sistem de reprezentanți} este o familie
$S=(x_i)_{i\in I}$ astfel încât
\[
 \begin{cases}
   x_i\not\sim x_j,&i\neq j,\\[2pt]
   \forall x\in M\ \exists i\in I:\ x\sim x_i.
 \end{cases}
\]
Pentru submulțimea corespunzătoare $S\subseteq M$, restricția
$\pi|_S:S\to M/{\sim}$ este bijectivă.

\medskip
\noindent\underline{Proprietatea de universalitate}

\noindent\textbf{Observație.} Fie $f:M\to X$. Definim
\[
 R_f:=\{(x,y)\in M\times M\mid f(x)=f(y)\}.
\]
Atunci $x\sim y\iff f(x)=f(y)$ definește relația de echivalență
indusă de $f$ pe $M$.

\noindent\fbox{T}\quad \textbf{Proprietatea de universalitate a mulțimii factor.}
Fie $\sim$ o relație de echivalență pe $M$ și fie
$\pi:M\to M/{\sim}$ proiecția canonică.
Pentru orice aplicație $f:M\to X$ cu $\sim\,\subseteq R_f$,
există o unică aplicație $\overline f:M/{\sim}\to X$ astfel încât
\[
 \overline f\circ\pi=f.
\]
\[
\begin{tikzcd}[column sep=large,row sep=large]
 M\arrow[r,"\pi"]\arrow[dr,"f"']
   & M/{\sim}\arrow[d,dashed,"{\overline f}"]\\
 & X
\end{tikzcd}
\]
% Manuscrisul se oprește după enunțul proprietății de universalitate.

\endgroup
